%% kapa: exam=1 round=20 year=2009 period=1 form=A subject=economics subject_ko=경제원론 seq=16 number=56
%% source: Q-Net 감정평가사 기출문제 (cst00309) artlSeq=5024178 file=extracted/1008/5024178/2009년 제20회 감정평가사 1차시험 기출문제/1교시.hwp
%% conversion: kordoc 4.24.51 -> tools/gen1.py (auto); review=proofread-sonnet; flags=none
\begin{question}{56}
절대적 구매력평가설에 관한 설명으로 옳지 \uline{않은} 것은?
\begin{choices}
  \item 장기적 환율결정이론이다.
  \item 구매력평가설이 성립하면 실질환율은 1이다.
  \item 명목환율은 두 나라 화폐 사이의 구매력 차이를 반영한다.
  \item 일물일가의 법칙에 근거하여 도출된다.
  \item 구매력평가설이 성립하면 순수출은 순자본유출과 크기가 같다.
\end{choices}
\end{question}
\answer{5}
